%%% FIELD proof-direct-support-transport-dual
Fix \(L\ge2\), a compatible marginal vector \(p\), and \(d\in\mathcal D\).  At the optimization level of this theorem the input is the full marginal vector \(p=(p_{rj})_{r,j}\).  If \(p\) is obtained from an observed randomized law, the observable map is \(P\mapsto(P(Y=j\mid A=r))_{r,j}\); the theorem itself uses only the supplied value of \(p\) and the maintained support \(\mathcal U_L\).  To keep the causal object distinct from its observed-law bound, define the scalar identified set induced by these marginals and the support restriction as
\[
 \Theta_{I,d}(p)=\{\theta_d(q):q\in\Gamma_L(p)\}.
\]
By \(\text{def:coupling-fiber}\), compatibility is exactly feasibility of
\[
 G_Lq=p,\qquad q\ge0.
 \tag{1}
\]
Summing the \(L\) equations for any fixed arm in (1) gives \(\sum_{u\in\mathcal U_L}q(u)=1\).  Hence \(\Gamma_L(p)\) is a closed subset of the finite probability simplex and is compact.  It is convex, and \(q\mapsto\theta_d(q)\) is linear by \(\text{def:unique-peak-indicators}\).  Therefore \(\Theta_{I,d}(p)\) is the compact interval
\[
 \Theta_{I,d}(p)=[\underline\theta_d(p),\overline\theta_d(p)].
 \tag{2}
\]
Thus any equality proved below is a sharp endpoint characterization, not merely an outer bound.

Put \(c^{(d)}_u=B_d(u)\).  The lower endpoint is the finite linear program
\[
 \min_{q\ge0}\left\{(c^{(d)})^{\mathsf T}q:G_Lq=p\right\}.
 \tag{3}
\]
It has an optimizer by compactness.  Its equality-form dual is
\[
 \max_{y\in\mathbb R^{3L}}
 \left\{p^{\mathsf T}y:G_L^{\mathsf T}y\le c^{(d)}\right\};
 \tag{4}
\]
the dual variable is unrestricted in sign because the primal constraints are equalities.  Applying the finite linear-programming strong-duality result \(\text{lem:finite-lp-strong-duality}\), in the exact form attributed there to Trevisan, Lecture 6, Theorem 6.5, shows that (4) attains the value of (3).  Since
\[
 (G_L^{\mathsf T}y)_u=\sum_{r\in\mathcal D}y_{r,u_r},
 \qquad
 p^{\mathsf T}y=\sum_{r\in\mathcal D}\sum_{j\in\mathcal Y_L}p_{rj}y_{rj},
\]
this is the asserted lower dual formula.  Applying the same theorem to the maximization primal
\[
 \max_{q\ge0}\left\{(c^{(d)})^{\mathsf T}q:G_Lq=p\right\}
\]
gives the attained dual minimum
\[
 \min_{z\in\mathbb R^{3L}}
 \left\{p^{\mathsf T}z:G_L^{\mathsf T}z\ge c^{(d)}\right\},
\]
which is the asserted upper formula.  Combining these equalities with (2) proves that the displayed dual values equal the sharp causal endpoints over all compatible latent couplings; that relationship was derived from the coupling fiber and was not imposed by definition.

We next derive the exact separator.  For \(r\in\{0,2\}\) and \(k\in\mathcal Y_L\), define
\[
 P_r^{\le}(k)=\max_{0\le j\le k}x_{rj},\qquad
 P_r^{<}(k)=\max_{0\le j<k}x_{rj},\qquad
 S_r^{>}(k)=\max_{k<j<L}x_{rj},
 \tag{5}
\]
where a maximum over an empty set is \(-\infty\).  The recursions
\[
 P_r^{\le}(k)=\max\{P_r^{\le}(k-1),x_{rk}\},
 \qquad
 S_r^{>}(k)=\max\{x_{r,k+1},S_r^{>}(k+1)\}
\]
with the evident boundary conventions compute all quantities in (5) in one forward and one backward pass, using \(O(L)\) exact additions and comparisons.

Fix \(u_1=k\).  By \(\text{def:non-valley-support}\), legality is equivalent to \(u_0\le k\) or \(u_2\le k\).  The legal pairs \((u_0,u_2)\) form the following five disjoint regions:
\[
\begin{array}{c|c|c}
\text{region}&\text{conditions}&B_d(u)\\ \hline
R_0&u_0>k,\ u_2\le k&\mathbf 1\{d=0\}\\
R_2&u_2>k,\ u_0\le k&\mathbf 1\{d=2\}\\
R_1&u_0<k,\ u_2<k&\mathbf 1\{d=1\}\\
R_{01}&u_0=k,\ u_2\le k&0\\
R_{12}&u_0<k,\ u_2=k&0.
\end{array}
 \tag{6}
\]
Indeed, if \(u_0>k\), legality forces \(u_2\le k\), giving \(R_0\); if \(u_2>k\), legality forces \(u_0\le k\), giving \(R_2\); otherwise both outer coordinates are at most \(k\), and the last three rows of (6) partition that case according as both inequalities are strict, only the first is an equality, or only the second is an equality.  On these regions, direct additive maximization gives
\[
\begin{aligned}
 V_{0d}(k)&=S_0^{>}(k)+x_{1k}+P_2^{\le}(k)-\sigma\mathbf 1\{d=0\},\\
 V_{2d}(k)&=P_0^{\le}(k)+x_{1k}+S_2^{>}(k)-\sigma\mathbf 1\{d=2\},\\
 V_{1d}(k)&=P_0^{<}(k)+x_{1k}+P_2^{<}(k)-\sigma\mathbf 1\{d=1\},\\
 V_{01,d}(k)&=x_{0k}+x_{1k}+P_2^{\le}(k),\\
 V_{12,d}(k)&=P_0^{<}(k)+x_{1k}+x_{2k}.
\end{aligned}
 \tag{7}
\]
Consequently,
\[
 \max_{u\in\mathcal U_L}
 \left\{\sum_{r\in\mathcal D}x_{r,u_r}-\sigma B_d(u)\right\}
 =\max_{k\in\mathcal Y_L}
 \max\{V_{0d}(k),V_{2d}(k),V_{1d}(k),V_{01,d}(k),V_{12,d}(k)\}.
 \tag{8}
\]
For the lower dual take \(x=y\) and \(\sigma=1\); its candidate is feasible exactly when (8) is at most zero.  For the upper dual take \(x=-z\) and \(\sigma=-1\), because the left side of (8) then equals
\[
 \max_{u\in\mathcal U_L}\left\{B_d(u)-\sum_{r\in\mathcal D}z_{r,u_r}\right\}.
\]
Keeping an argmax in every prefix and suffix recursion also returns a maximizing legal triple, and hence a violated dual constraint whenever the maximum is positive.  This proves the \(O(L)\) exact separation claim.  Because \(\mathcal U_L\) is finite, adding a newly violated constraint terminates after at most \(|\mathcal U_L|\) additions; equivalently, one may enumerate all legal columns and solve the resulting finite rational program.  Thus rational input gives a fully exact finite computation without any catalogue of canonical facets.

For completeness, we prove the stronger rational-input complexity and sparsity corollary.  A triple is excluded precisely when, for \(u_1=k\), both outer coordinates lie in \(\{k+1,\ldots,L-1\}\).  Hence
\[
 |\mathcal U_L|
 =L^3-\sum_{k=0}^{L-2}(L-1-k)^2
 =L^3-\sum_{m=1}^{L-1}m^2
 =L^3-\frac{(L-1)L(2L-1)}6
 =O(L^3).
 \tag{9}
\]
The explicit feasibility program and each of the six primal endpoint programs therefore have a number of binary-encoded rational coefficients polynomial in \(L\) and in the encoding length of \(p\).  The rational linear-programming algorithm in \(\text{lem:rational-lp-polynomial-time}\) decides compatibility and, when compatible, returns the exact optimum of every one of these six bounded programs in polynomial time.

To expose the optimizer construction and its support bound, replace the redundant full marginal equations by the matrix \(M_L\) whose rows are one normalization row and the \(3(L-1)\) rows indexed by \((r,j)\in\mathcal D\times\{1,\ldots,L-1\}\):
\[
 (M_L)_{\varnothing,u}=1,
 \qquad
 (M_L)_{(r,j),u}=\mathbf 1\{u_r=j\}.
 \tag{10}
\]
For \(b(p)=(1,(p_{rj})_{r,1\le j<L})^{\mathsf T}\), the simplex identities imply
\[
 \Gamma_L(p)=\{q\ge0:M_Lq=b(p)\}.
 \tag{11}
\]
Indeed, the omitted level-zero equation for arm \(r\) is
\[
 \sum_{u:u_r=0}q(u)
 =1-\sum_{j=1}^{L-1}p_{rj}=p_{r0}.
\]
The matrix in (10) has only \(3L-2\) rows.  An optimal basic feasible solution consequently has at most \(3L-2\) positive entries.  Such a solution is obtainable within the stated polynomial bound.  After computing the rational optimal value \(v\), start with the nonempty compact rational polytope obtained from (11) by adding the equality \((c^{(d)})^{\mathsf T}q=v\).  In a fixed ordering of its \(O(L^3)\) coordinates, sequentially minimize each coordinate over the face left by all preceding minimizations and add the equality fixing that coordinate at its attained minimum.  Each minimum is attained at a vertex of the original optimal face, so Cramer's rule applied to a submatrix of the original rational system bounds its encoding length polynomially in the original input.  Each step is therefore an exact rational linear program of polynomial encoding length, and there are \(O(L^3)\) steps.  At the end every coordinate is fixed, so the remaining face is a singleton.  A singleton face of the optimal face is an extreme point of (11), hence a basic optimal solution.  If the positive columns of \(M_L\) at that point were dependent, a sufficiently small two-sided null perturbation would express it as the midpoint of two distinct feasible points, contradicting extremality.  Thus its support has size at most \(\operatorname{rank}(M_L)\le3L-2\).

The same support bound also has an explicit support-elimination certificate.  If a rational optimizer \(q\) has support \(S\) with \(|S|>3L-2\), the columns \((M_{L,u})_{u\in S}\) are dependent, so exact Gaussian elimination returns nonzero rational \(a\), supported on \(S\), with \(M_La=0\).  The normalization row implies \(\sum_ua(u)=0\), so \(a\) has both signs.  For sufficiently small \(\varepsilon>0\), both \(q+\varepsilon a\) and \(q-\varepsilon a\) are feasible.  Optimality for either the minimum or maximum of the linear cost forces
\[
 \sum_{u\in\mathcal U_L}B_d(u)a(u)=0;
\]
otherwise one of the two signs would improve the objective.  Therefore
\[
 t=\min_{u\in S:a(u)<0}\frac{q(u)}{-a(u)}>0
\]
makes \(q+ta\) a rational optimizer with strictly smaller support.  At most \(|\mathcal U_L|-(3L-2)\) repetitions give a rational endpoint optimizer supported on at most \(3L-2\) legal triples.  The preceding sequential-face construction supplies such a basic optimizer with a polynomial number of exact rational-LP calls, while this elimination argument certifies the support bound equation by equation.  Thus compatibility, all six endpoints, and sparse rational endpoint-attaining mixtures are exactly computable in time polynomial in \(L\) and the input bit length, while the all-\(L\) canonical facet stream remains unnecessary for scalar endpoint evaluation.
%%% FIELD statement-rational-lp-polynomial-time
Every finite linear program with binary-encoded rational coefficients can be solved by an exact algorithm whose running time is polynomial in the total input encoding length: the algorithm decides feasibility and, when the feasible objective is bounded, returns an exact rational optimal value and an exact rational optimal solution.
%%% FIELD statement-rational-endpoint-computation
Fix \(L\ge2\).  Given a rational vector \(p\in\prod_{d\in\mathcal D}\Delta^{L-1}\) in binary encoding, whether \(\Gamma_L(p)\) is empty is decidable in time polynomial in \(L\) and the total input bit length.  If it is nonempty, all six values \(\underline\theta_d(p)\) and \(\overline\theta_d(p)\), \(d\in\mathcal D\), are exactly computable within the same complexity class.  For each endpoint, the algorithm also constructs a rational optimizer \(q^\star\in\Gamma_L(p)\) with \(|\operatorname{supp}(q^\star)|\le3L-2\).  These conclusions use \(|\mathcal U_L|=L^3-(L-1)L(2L-1)/6=O(L^3)\) and do not require enumeration of \(\mathcal H_L\).
%%% FIELD proof-rational-endpoint-computation
For a fixed middle response \(u_1=k\), the illegal strict valleys have independently chosen outer coordinates \(u_0,u_2\in\{k+1,\ldots,L-1\}\), and therefore number \((L-1-k)^2\).  Summing over \(k=0,\ldots,L-2\) gives
\[
 |\mathcal U_L|
 =L^3-\sum_{k=0}^{L-2}(L-1-k)^2
 =L^3-\frac{(L-1)L(2L-1)}6.
 \tag{12}
\]
In particular, the legal-incidence matrix in \(\text{thm:direct-support-transport-dual}\) has \(3L\) rows and \(O(L^3)\) columns, all with zero--one entries.  For rational \(p\), its binary description has length polynomial in \(L\) and the input encoding length.  Apply \(\text{lem:rational-lp-polynomial-time}\) first to the feasibility system \(G_Lq=p\), \(q\ge0\).  This decides whether \(\Gamma_L(p)\) is empty.  When it is nonempty, apply the same result to the minimization and maximization programs with cost \((B_d(u))_{u\in\mathcal U_L}\), for each of the three values of \(d\).  Their feasible set is contained in the probability simplex and is therefore bounded.  The direct-support dual theorem identifies their six exact optimal values with \(\underline\theta_d(p)\) and \(\overline\theta_d(p)\).  A constant number of polynomial-time calls remains polynomial, proving the computation claim.

It remains to construct the stated optimizer.  Let \(M_L\) have the normalization row and the \(3(L-1)\) nonzero-level marginal rows,
\[
 (M_L)_{\varnothing,u}=1,
 \qquad
 (M_L)_{(r,j),u}=\mathbf 1\{u_r=j\},
 \qquad 1\le j<L.
\]
For \(b(p)=(1,(p_{rj})_{r,1\le j<L})^{\mathsf T}\), the simplex equations give
\[
 \Gamma_L(p)=\{q\ge0:M_Lq=b(p)\},
 \tag{13}
\]
because the omitted equation equals \(1-\sum_{j=1}^{L-1}p_{rj}=p_{r0}\) for every arm \(r\).  For a polynomial-time basic-solution construction, first compute the rational optimal value \(v\).  Add \(c^{\mathsf T}q=v\) to (13), fix an ordering of the \(O(L^3)\) coordinates, and sequentially minimize each coordinate over the nonempty compact face left by preceding minimizations, fixing that coordinate at its exact attained value after each call.  Each minimum is attained at a vertex of the original optimal face; Cramer's rule for a nonsingular submatrix of the original rational system therefore gives it encoding length polynomial in the original input.  The cited algorithm consequently solves every one of these polynomial-size rational programs in polynomial time.  After all coordinates have been fixed, the remaining face is a singleton, hence an extreme point of (13).  Its positive columns of \(M_L\) are linearly independent: otherwise a sufficiently small two-sided null perturbation would express it as the midpoint of two distinct feasible points.  Consequently this rational optimizer has support at most \(\operatorname{rank}(M_L)\le3L-2\).  The number of calls is polynomial because there are \(O(L^3)\) coordinates.

For an equation-by-equation support certificate starting from any rational optimizer \(q\), suppose its support \(S\) has more than \(3L-2\) elements.  Then the corresponding columns of the \((3L-2)\)-row matrix \(M_L\) are dependent.  Exact Gaussian elimination produces nonzero rational \(a\), supported on \(S\), with \(M_La=0\).  The normalization row makes \(a\) have both signs.  Both \(q+\varepsilon a\) and \(q-\varepsilon a\) are feasible for all sufficiently small positive \(\varepsilon\).  If \(c_u=B_d(u)\) is the endpoint cost, optimality of \(q\) under these two perturbations implies \(c^{\mathsf T}a=0\).  Thus
\[
 t=\min_{u\in S:a(u)<0}\frac{q(u)}{-a(u)}
\]
is positive and rational, and \(q+ta\) is a rational optimizer with at least one fewer positive coordinate.  Repeating yields a rational optimizer with the asserted support.  The sequential-face construction above establishes the polynomial-time guarantee without requiring a bound on the intermediate fractions along this optional elimination path.  Only the legal incidence matrix, the endpoint cost, and the observed marginal vector enter either construction, so no enumeration of \(\mathcal H_L\) is used.
%%% FIELD statement-direct-support-postimage
Fix \(L\ge2\).  Let \(G_L\) be the \(3L\)-by-\(|\mathcal U_L|\) legal-response incidence matrix with rows indexed by \((r,j)\in\mathcal D\times\mathcal Y_L\), columns indexed by \(u\in\mathcal U_L\), and entries \((G_L)_{(r,j),u}=\mathbf 1\{u_r=j\}\).  For \(d\in\mathcal D\), let \(c^{(d)}_u=B_d(u)\), and regard \(p=(p_{rj})_{r,j}\) here as the full \(3L\)-vector of arm probabilities.  For every compatible \(p\),
\[
 \underline\theta_d(p)=
 \max_{y\in\mathbb R^{3L}}
 \left\{\sum_{r\in\mathcal D}\sum_{j\in\mathcal Y_L}p_{rj}y_{rj}:
 \sum_{r\in\mathcal D}y_{r,u_r}\le B_d(u)\quad\forall u\in\mathcal U_L\right\},
\]
and
\[
 \overline\theta_d(p)=
 \min_{z\in\mathbb R^{3L}}
 \left\{\sum_{r\in\mathcal D}\sum_{j\in\mathcal Y_L}p_{rj}z_{rj}:
 \sum_{r\in\mathcal D}z_{r,u_r}\ge B_d(u)\quad\forall u\in\mathcal U_L\right\}.
\]
Both dual optima are attained.  Moreover, feasibility of either dual candidate has an \(O(L)\) exact separation oracle.  For arrays \(x=(x_{rj})\) and \(\sigma\in\{-1,1\}\), all prefix and suffix maxima can be computed in two passes, after which
\[
 \max_{u\in\mathcal U_L}\left\{\sum_{r\in\mathcal D}x_{r,u_r}-\sigma B_d(u)\right\}
\]
is the maximum, over \(k\in\mathcal Y_L\), of five explicitly evaluable region scores corresponding to \(u_0>k,u_2\le k\); \(u_2>k,u_0\le k\); \(u_0<k,u_2<k\); \(u_0=k,u_2\le k\); and \(u_0<k,u_2=k\).  Thus scalar endpoints can be computed by a finite linear program with this separator, using exact rational arithmetic whenever \(p\) is rational.  The complete all-\(L\) canonical facet stream is needed to catalogue the joint polytope, but is not needed merely to evaluate a scalar endpoint.
